\section{The modular cubic boundary in characteristic three}\label{sec:higher-three-modular}
This section identifies the nonsplit node module that forces the additional unipotent transformations in characteristic three. Assume $n\ge3$, $d\ge n+3$, and put $m=d-3$.

The cubic-boundary calculation of Section~\ref{sec:cubic-cech} gives
\begin{equation}\label{eq:three-extension}
0\longrightarrow P\longrightarrow V\longrightarrow
E=\bigoplus_{|I|=n-2}E_I\longrightarrow0,
\qquad
0\longrightarrow P_0\longrightarrow P\longrightarrow
\bigoplus_{|J|=n-1}Q_J\longrightarrow0.
\end{equation}
Here $P_0$ is constant of dimension $\binom mn$ and $Q_J=\mathbf F_3[Z_J]/\mathbf F_3\mathbf1$. We do not split the second sequence. In characteristic three its failure to split is essential.

For clarity, the first sequence follows directly from the intersection resolution. A cubic of dimension at least two has no positive structure-sheaf cohomology; a plane cubic has one-dimensional $H^1$; and a line section has length three. These statements follow from the hypersurface sequence and the projective-space calculation in \citet[Lemmas~30.8.1--30.8.2]{StacksProjective}. Thus the only positive row is the elliptic row at bidegree $(n-2,1)$. The resolution ends in degree $n-1$, so no differential can leave that row toward degree $n$ in the zero row. Replacing the three point values by one diagonal constant gives the constant skeleton complex, whose top cohomology is $P_0$. The quotient is supported only in the last degree, giving the second sequence in~\eqref{eq:three-extension}. In particular the entire action on $P$ is determined by actual permutations of the three-point sets, since all other zero-row terms and maps are constant.

Allow the cubic to become three labelled hyperplanes $a,b,c$, transverse to the other $m$ hyperplanes. The plane sections are ordinary triangles. The grouped intersection resolution remains exact: locally its cubic equation is a product of coordinate variables disjoint from those defining the remaining hyperplanes, and the alternating ideal complex is exact monomial by monomial. Refining the grouped resolution gives the fully labelled one. The hypersurface calculation keeps all ranks constant across this specialization and makes the cohomology and intersection maps commute with base change. The ordinary coefficient open containing the triangle and the smooth-cubic open is normal and integral. Section~\ref{sec:cubic-generic} therefore identifies the finite image at this common ordinary fibre.

Let
\[
U_0=\mathbf F_3^d,\qquad u=\sum_l e_l,\qquad
W=U_0/\mathbf F_3u.
\]
The labelled arrangement complex has differential $u\wedge$, and hence
\begin{equation}\label{eq:three-arrangement}
V=\bigwedge^nU_0/(u\wedge\bigwedge^{n-1}U_0)=\bigwedge^nW.
\end{equation}
An $n$-fold point cochain maps to the exterior product of its labels. The parameter space of ordered general-position hyperplanes is irreducible; quotienting the three labels $a,b,c$ gives a connected $S_3$ cover. Their physical permutations are therefore actual elements of the grouped boundary image, acting on~\eqref{eq:three-arrangement} by exterior powers.

Put $A=e_a$, $B=e_b$, let $B_0$ be spanned by the $m$ outside vectors, and put $S=\sum_{j=1}^m e_j$. Then $e_c=-A-B-S$ and $W=\mathbf F_3A\oplus\mathbf F_3B\oplus B_0$. Refinement sends the three values at $Z_J$ to $A\wedge e_J,B\wedge e_J,e_c\wedge e_J$ and the hyperplane-only points to their outside cochains. Thus
\begin{align}
P={}&\bigwedge^nB_0\oplus
 A\wedge\bigwedge^{n-1}B_0\oplus
 B\wedge\bigwedge^{n-1}B_0,\label{eq:three-P}\\
P_0={}&\bigwedge^nB_0,
\qquad E=\bigoplus_I\mathbf F_3\overline w_I,
\qquad w_I=A\wedge B\wedge e_I.\nonumber
\end{align}
The outside $n$-fold cochains together with $A\wedge e_J$ and $B\wedge e_J$ form the displayed independent basis; the $e_c$ values satisfy $e_c=-A-B-S$. Its dimension agrees with~\eqref{eq:three-extension}, so these identifications are exact. Let $\phi_I$ denote the coordinate row of $V\to E_I$ in this quotient basis.

At an individual node index $J$, the transposition $(a\ b)$ swaps $A\wedge e_J$ and $B\wedge e_J$, fixes all other such pairs, and fixes $P_0$. The action on each $Q_J$ is faithful. Theorem~\ref{thm:joint} therefore identifies the quotient of the actual boundary image acting on $P$ and $E$ with
\begin{equation}\label{eq:three-graded-group}
L=\prod_J S_3\times\prod_I\{1,-1\}.
\end{equation}
This specifies the actual modular node action, not just its composition factors.

\section{Separating the elliptic characters in characteristic three}\label{sec:higher-three-kernel}
This section splits the finite group between $P$ and $E$, while preserving the nonsplit node module inside $P$.

Let $G_B$ be the boundary image and let $N$ be the kernel of its action on $P$ and $E$. Its elements have the form $1+b$ with $b:E\to P$, so $N$ is an additive $\mathbf F_3$ subspace of $\operatorname{Hom}(E,P)$. Lift the element of~\eqref{eq:three-graded-group} acting trivially on $P$ and by $-1$ on every $E_I$. Every such lift has block form $\left(\begin{smallmatrix}1&b\\0&-1\end{smallmatrix}\right)$ and square one. Use its negative eigenspace as a complement to $P$, so that the lift $c_0$ is $\operatorname{diag}(1,-1)$.

For $x=\left(\begin{smallmatrix}A_0&b\\0&D\end{smallmatrix}\right)\in G_B$, direct multiplication gives
\[
[x,c_0]=1+2bD^{-1}\in N.
\]
Thus $bD^{-1}\in N$, and multiplying $x$ on the left by $1-bD^{-1}$ removes its off-diagonal block. The centralizer of $c_0$ maps onto $L$ and meets $N$ trivially. Consequently $G_B$ contains a block-diagonal copy of $L$, and $N$ is stable under $b\mapsto A_0bD^{-1}$.

The independent sign on $E_I$ projects $N$ to its $I$-th column: subtract the original element from its conjugate, which multiplies that column by $-2=1$. Put
\[
N_I=\{v\in P:v\phi_I\in N\}.
\]
Each $N_I$ is stable under every independent node permutation. None of this asserts a splitting of $P_0\subset P$.

\section{Modular norms force actual transvections}\label{sec:higher-three-force}
This section obtains supported rank-one transformations from a physical three-cycle and identifies their rows as arrangement cochains.

Let $\sigma$ be the physical cycle $a\mapsto b\mapsto c\mapsto a$. It acts by $+1$ on every triangle quotient $E_I$. In the chosen complement, write its lift of $\overline w_I$ as $x_I=w_I+p_I$, $p_I\in P$. The off-diagonal column of $\sigma$ is $b_I=(\sigma-1)x_I$. Section~\ref{sec:higher-three-kernel} gives $b_I\in N_I$, and stability gives
\[
z_I=(\sigma-1)b_I=(\sigma-1)^2x_I=(1+\sigma+\sigma^2)x_I\in N_I.
\]
The norm sends $P$ into $P_0$: it is zero on $P_0$ and sends $A\wedge e_J$ and $B\wedge e_J$ to $-S\wedge e_J$. Moreover
\[
\sigma w_I=w_I-B\wedge S\wedge e_I,
\qquad \sigma^2w_I=w_I+A\wedge S\wedge e_I.
\]
Therefore
\begin{equation}\label{eq:three-norm}
z_I=(A-B)\wedge S\wedge e_I+q_I,
\qquad q_I\in P_0.
\end{equation}
For $j\notin I$ and $J=I\cup\{j\}$, apply $t_J-1$, where $t_J$ is the independent node transposition $(a\ b)$ at $J$. It kills all other node components and $P_0$, and multiplies the displayed $J$ component by $-2=1$. Hence $(A-B)\wedge e_J\in N_I$. The actual image contains
\begin{equation}\label{eq:three-supported}
1+\bigl[(e_a-e_b)\wedge e_I\wedge e_j\bigr]\phi_I.
\end{equation}
Its centre lies in $P$ and its row vanishes there, so this is a transvection on the full representation. Regrouping any three of the labelled hyperplanes as the cubic gives the same conclusion for every allowed choice of labels at the fixed fibre~\eqref{eq:three-arrangement}.

For an ordered $(n+1)$-element label set $U$, let $\delta_U$ have the alternating boundary values on its $n$-faces and value zero on the other coordinate $n$-products. It annihilates $u\wedge e_K$, since the two possible boundary terms cancel. Thus it is a row on $\bigwedge^nW$. For $U=\{a,b,c\}\cup I$, it vanishes on~\eqref{eq:three-P}; on $w_{I'}$ it is zero unless $I'=I$, and then is $1$ or $-1$. Accordingly
\begin{equation}\label{eq:three-row}
\phi_I=\pm\delta_{\{a,b,c\}\cup I}.
\end{equation}
The sign is absorbed by taking the inverse transvection. This quotient-row calculation is independent of the complement used to exhibit $N$.

\section{Elementary generation in characteristic three}\label{sec:higher-three-group}
This section converts the supported transformations into all elementary matrices, completing the characteristic-three image calculation without a classification theorem for finite linear groups.

Choose $e_1,\ldots,e_{d-1}$ as a basis of $W$, so $e_d=-\sum_{r<d}e_r$. Coordinate vectors of $V$ are $e_A$ for $n$-subsets $A$ of the first $d-1$ labels. Fix $R=\{r_0,\ldots,r_n\}$ in increasing order, and write
\[
f_i=(-1)^i e_{R\setminus\{r_i\}}.
\]
For distinct $i,j,k$, apply~\eqref{eq:three-supported} using $a=r_j$, $b=r_k$, $c=d$, outside set $I=R\setminus\{r_i,r_j,r_k\}$, and the additional label $r_i$. Its row in~\eqref{eq:three-row} is a nonzero multiple of $f_i^*$ on the entire coordinate space: only the face $R\setminus\{r_i\}$ avoids $d$. Its centre is a common nonzero multiple of $f_j+f_k$, by the alternating face signs. Taking an inverse if necessary gives every
\begin{equation}\label{eq:three-triple}
1+(f_j+f_k)f_i^*.
\end{equation}
These transformations fix all coordinates outside this clique of $n+1$ faces.

Since $n+1\ge4$, for each $i\ne j$ choose two other distinct indices $k,l$. The same-row transformations commute, and
\begin{align*}
&[1+(f_j+f_k)f_i^*]\,[1+(f_j+f_l)f_i^*]\,[1-(f_k+f_l)f_i^*]\\
&\hspace{35mm}=1+2f_jf_i^*.
\end{align*}
Squaring gives $1+f_jf_i^*$. Thus the image contains all elementary transformations within each clique.

The graph of coordinate $n$-subsets joined by one-element replacements is connected. Every edge belongs to the clique of faces of its $(n+1)$-element union, so both elementary directions occur along every edge. For distinct coordinate vertices $A,B,C$,
\[
[1+E_{CB},1+E_{BA}]=1+E_{CA}.
\]
Commuting along a simple path gives every off-diagonal elementary matrix. They generate $\mathrm{SL}_h(\mathbf F_3)$ by Gaussian elimination; the determinant-one diagonal operations follow from the usual two-by-two elementary identities.

Finally the block-diagonal complement of Section~\ref{sec:higher-three-kernel} contains the element acting by $-1$ on one $E_I$ and identically on $P$ and all other elliptic lines. Its determinant is $-1$. The image therefore contains $\mathrm{SL}_h(\mathbf F_3)$ and both determinants, hence equals $\mathrm{GL}_h(\mathbf F_3)$. The geometric-to-arithmetic containment proves the same equality over every perfect ground field. Together with the first-degree chart in Section~\ref{sec:three-seed}, this completes the characteristic-three case of Theorem~\ref{thm:higher-images}.
